\section{Teaching Activities}\label{teaching-activities}

\subsection{Courses}\label{courses-1}

In this section I list all the official curricular units in which I have
been involved at IST. They are grouped in PhD level courses, MSc level
courses, BSc level courses, and Tuturial Courses. The former three are
conventional courses with group lectures and practice/labs. The latter
is of tutorial nature with meetings with students (individually or in
group) to guide their preparation of the MSc thesis projects. For the
former two cases, I include, whenever available, the ranks given by
students to my teaching performance, according to the student's surveys
organized by IST pedagogical council program QUC (IST Course Unit
Quality System).

\subsubsection{PhD Level Courses (3rd
Cycle)}\label{phd-level-courses-3rd-cycle}

\begin{enumerate}
\def\labelenumi{\arabic{enumi}.}
\item
  \ul{Estimation and Classification}
\end{enumerate}

2025/2026, 1\textsuperscript{st} Semester, PDEEC, Teaching

2024/2025, 1\textsuperscript{st} Semester, PDEEC, Teaching

\subsubsection{\texorpdfstring{MSc Level Courses (2\textsuperscript{nd}
Cycle)}{MSc Level Courses (2nd Cycle)}}\label{msc-level-courses-2nd-cycle}

\begin{enumerate}
\def\labelenumi{\arabic{enumi}.}
\setcounter{enumi}{1}
\item
  \ul{Distributed Real Time Control Systems}
\end{enumerate}

2025/2026, 2\textsuperscript{nd} Semester, MEEC21, Coordination and
Teaching

Students Rating: NA (lectures and labs)

2024/2025, 2\textsuperscript{nd} Semester, MEEC21, Coordination and
Teaching

Students Rating: 8.75/9 (lectures), 9/9 (labs)

2023/2024, 2\textsuperscript{nd} Semester, MEEC21, Coordination and
Teaching

Students Rating: 8.62/9 (lectures), 8.5/9 (labs)

2022/2023, 2\textsuperscript{nd} Semester, MEEC21, Coordination and
Teaching

Students Rating: NA (lectures and labs)

\begin{quote}
2021/2022, 2\textsuperscript{nd} Semester, MEEC21, Coordination and
Teaching

Students Rating: 8.69/9 (lectures and labs)

2020/2021, 1\textsuperscript{st} Semester, MEAer, MEEC, Coordination and
Teaching

Students Rating: 8.75/9 (lectures), 9/9 (labs)

2019/2020, 1\textsuperscript{st} Semester, MEAer, MEEC, Coordination and
Teaching

Students Rating: 8.12/9 (lectures), 8.38/9 (labs)

2018/2019, 1\textsuperscript{st} Semester, MEAer, MEEC, Coordination and
Teaching

Students Rating: 8.12/9 (lectures), 8.62 (labs)

2017/2018, 1\textsuperscript{st} Semester, MEAer, MEEC, Coordination and
Teaching

Students Rating: 7.25/9 (lectures), 7.88 (labs)

2016/2017, 1\textsuperscript{st} Semester, MEAer, MEEC, Coordination and
Teaching

Students Rating: 8.88/9 (lectures), 9/9 (labs)

2015/2016, 1\textsuperscript{st} Semester, MEAer, MEEC, Coordination and
Teaching

Students Rating: 6.88/9 (lectures)

2013/2014, 1\textsuperscript{st} Semester, MEAer, MEEC, Coordination and
Teaching

Students Rating: 8.25/9 (lectures)

2012/2013, 1\textsuperscript{st} Semester, MEAer, MEEC, Coordination and
Teaching

Students Rating: 8.12/9 (lectures)

2009/2010, 1\textsuperscript{st} Semester, MEAer, MEEC, Coordination and
Teaching

Students Rating: 7.38/9 (lectures)

2008/2009, 1\textsuperscript{st} Semester, MEAer, MEEC, Coordination and
Teaching

Students Rating: NA (lectures)

2007/2008, 1\textsuperscript{st} Semester, MEAer, MEEC, Coordination and
Teaching

Students Rating: NA (lectures)
\end{quote}

\begin{enumerate}
\def\labelenumi{\arabic{enumi}.}
\setcounter{enumi}{2}
\item
  \ul{Modelling, Identification, and Digital Control}
\end{enumerate}

2013/2014, 1\textsuperscript{st} Semester, MEEC, Teaching

Students Rating: 7.62/9 (labs)

2012/2013, 1\textsuperscript{st} Semester, MEEC, Teaching

Students Rating: 7.5/9 (labs)

2011/2012, 1\textsuperscript{st} Semester, MEEC, Teaching

Students Rating: 8.56/9 (labs)

2010/2011, 1\textsuperscript{st} Semester, MEEC, Coordination and
Teaching

Students Rating: 6.6/9 (lectures), 6.6/9 (labs)

2006/2007, 1\textsuperscript{st} Semester, LEEC-pB, Coordination and
Teaching

Students Rating: 3.83/5 (lectures)

2005/2006, 1\textsuperscript{st} Semester, LEEC-pB, MEEC-pB,
Coordination and Teaching

Students Rating: 3.91/5 (Lectures), 4.04/5 (Labs)

2004/2005, 1\textsuperscript{st} Semester, LEEC-pB, Coordination and
Teaching

Students Rating: 4.05/5 (Lectures), 4.09/5 (Labs)

2003/2004, 1\textsuperscript{st} Semester, LEEC-pB, Coordination and
Teaching

Students Rating: NA (lectures, labs)

2002/2003, 1\textsuperscript{st} Semester, LEEC-pB, Teaching

Students Rating: NA (labs)

\begin{enumerate}
\def\labelenumi{\arabic{enumi}.}
\setcounter{enumi}{3}
\item
  \ul{Machine Learning}
\end{enumerate}

2025/2026, 1\textsuperscript{st} Semester, Multiple courses,
Coordination and Teaching

Students Rating: NA (lectures)

2024/2025, 1\textsuperscript{st} Semester, Multiple courses,
Coordination and Teaching

Students Rating: 7.88/9 (lectures)

2023/2024, 1\textsuperscript{st} Semester, Multiple courses,
Coordination and Teaching

Students Rating: 7.88/9 (lectures)

2009/2010, 1\textsuperscript{st} Semester, MEBiom, MEEC, Coordination
and Teaching

Students Rating: 5.06/9 (lectures)

2008/2009, 1\textsuperscript{st} Semester, MEBiom, MEEC, Teaching

Students Rating: 6.94/9 (labs)

\begin{enumerate}
\def\labelenumi{\arabic{enumi}.}
\setcounter{enumi}{4}
\item
  \ul{Autonomous Systems}
\end{enumerate}

2021/2022, 2\textsuperscript{nd} Semester, LEEC21, MEAer21, MEEC21,
Min-IA, Teaching

Students Rating: 8.75/9 (lectures and labs)

\begin{enumerate}
\def\labelenumi{\arabic{enumi}.}
\setcounter{enumi}{5}
\item
  \ul{Computer Controlled Systems}
\end{enumerate}

2014/2015, 1\textsuperscript{st} Semester, MEAer, MEEC, Teaching

Students Rating: 8.07/9 (labs)

\begin{enumerate}
\def\labelenumi{\arabic{enumi}.}
\setcounter{enumi}{6}
\item
  \ul{Optimal and Adaptive Control}
\end{enumerate}

2012/2013, 1\textsuperscript{st} Semester, MEAer, Teaching

Students Rating: 8.81/9 (labs)

2011/2012, 1\textsuperscript{st} Semester, MEAer, Teaching

Students Rating: 7.75/9 (labs)

2010/2011, 1\textsuperscript{st} Semester, MEAer, Teaching

Students Rating: NA (labs)

2004/2005, 2\textsuperscript{nd} Semester, LEA-pB, Teaching

Students Rating: NA (labs)

\begin{enumerate}
\def\labelenumi{\arabic{enumi}.}
\setcounter{enumi}{7}
\item
  \ul{Robotics}
\end{enumerate}

2011/2012, 2\textsuperscript{nd} Semester, MEEC, MEIC-A, MEBiom,
Teaching

Students Rating: NA (labs)

1996/1997, 1\textsuperscript{st} Semester, LEIC-pB, Teaching

Students Rating: 3.70/5 (labs)

\begin{enumerate}
\def\labelenumi{\arabic{enumi}.}
\setcounter{enumi}{8}
\item
  \ul{State-Space Control}
\end{enumerate}

2006/2007, 2\textsuperscript{nd} Semester, LEEC-pB, Teaching

Students Rating: 4.02/5 (practice), 4.06/5 (labs)

2005/2006, 2\textsuperscript{nd} Semester, LEEC-pB, Teaching

Students Rating: 4.10/5 (practice), 4.03/5 (labs)

2004/2005, 2\textsuperscript{nd} Semester, LEEC-pB, Teaching

Students Rating: 3.58/5 (labs)

\begin{enumerate}
\def\labelenumi{\arabic{enumi}.}
\setcounter{enumi}{9}
\item
  \ul{Automation and Industrial Processes}
\end{enumerate}

2002/2003, 2\textsuperscript{nd} Semester, LEEC-pB, Teaching

Students Rating: 4.07/5 (labs)

\begin{enumerate}
\def\labelenumi{\arabic{enumi}.}
\setcounter{enumi}{10}
\item
  \ul{Computer Vision}
\end{enumerate}

2001/2002, 2\textsuperscript{nd} Semester, LEIC-pB, Teaching

Students Rating: NA (lectures, labs)

1997/1998, 2\textsuperscript{nd} Semester, LEIC-pB, Teaching

Students Rating: NA (labs)

\begin{enumerate}
\def\labelenumi{\arabic{enumi}.}
\setcounter{enumi}{11}
\item
  \ul{Neural Networks}
\end{enumerate}

1997/1998, 2\textsuperscript{nd} Semester, LEEC-pB, Teaching

Students Rating: NA (labs)

1996/1997, 2\textsuperscript{nd} Semester, LEEC-pB, Teaching

Students Rating: 4.43/5 (labs)

1995/1996, 2\textsuperscript{nd} Semester, LEEC-pB, Teaching

Students Rating: 4.43/5 (labs)

\subsubsection{\texorpdfstring{BSc Level Courses (1\textsuperscript{st}
Cycle)}{BSc Level Courses (1st Cycle)}}\label{bsc-level-courses-1st-cycle}

\begin{enumerate}
\def\labelenumi{\arabic{enumi}.}
\setcounter{enumi}{12}
\item
  \ul{Modeling and Simulation}
\end{enumerate}

2023/2024, 1\textsuperscript{st} Semester, LEEC, Teaching

Students Rating: 6.25/9 (lectures)

2022/2023, 1\textsuperscript{st} Semester, LEEC, Teaching

Students Rating: NA (lectures), 7.62/9 (labs)

2019/2020, 2\textsuperscript{nd} Semester, MEEC, Teaching

Students Rating: 8.12/9 (labs)

2018/2019, 2\textsuperscript{nd} Semester, MEEC, Teaching

Students Rating: 7.25/9 (labs)

2017/2018, 2\textsuperscript{nd} Semester, MEEC, Teaching

Students Rating: 6.81/9 (labs)

2010/2011, 2\textsuperscript{nd} Semester, MEEC, Teaching

Students Rating: 8.12/9 (labs)

2009/2010, 2\textsuperscript{nd} Semester, MEEC, Teaching

Students Rating: 7.28/9 (labs)

\begin{enumerate}
\def\labelenumi{\arabic{enumi}.}
\setcounter{enumi}{13}
\item
  \ul{Signals and Systems}
\end{enumerate}

\begin{quote}
2019/2020, 1\textsuperscript{st} Semester, MEBiom, MEEC, Teaching

Students Rating: 6.88/9 (practice)

2008/2009, 2\textsuperscript{nd} Semester, MEEC, Teaching

Students Rating: 6.9/9 (practice)

2003/2004, 2\textsuperscript{nd} Semester, LEEC-pB, Teaching

Students Rating: 4.01/5 (practice)

2001/2002, 2\textsuperscript{nd} Semester, LEEC-pB, Teaching

Students Rating: NA (practice)
\end{quote}

\begin{enumerate}
\def\labelenumi{\arabic{enumi}.}
\setcounter{enumi}{14}
\item
  \ul{Portfolio MEEC}
\end{enumerate}

\begin{quote}
2018/2019, 1\textsuperscript{st} Semester, MEEC, Teaching

Students Rating: 8.25/9 (labs)

2017/2018, 1\textsuperscript{st} Semester, MEEC, Teaching

Students Rating: 7.75/9 (labs)
\end{quote}

\begin{enumerate}
\def\labelenumi{\arabic{enumi}.}
\setcounter{enumi}{15}
\item
  \ul{Algorithms and Data Structures}
\end{enumerate}

2017/2018, 2\textsuperscript{nd} Semester, MEEC, Teaching

Students Rating: 8.12/9 (labs)

\begin{enumerate}
\def\labelenumi{\arabic{enumi}.}
\setcounter{enumi}{16}
\item
  \ul{Introduction to Control}
\end{enumerate}

2014/2015, 1\textsuperscript{st} Semester, MEAer, Teaching

Students Rating: 6.5/9 (practice)

\begin{enumerate}
\def\labelenumi{\arabic{enumi}.}
\setcounter{enumi}{17}
\item
  \ul{Control}
\end{enumerate}

\begin{quote}
2003/2004, 2\textsuperscript{nd} Semester, LEA-pB, LEN-pB, Teaching

Students Rating: NA (labs)

2002/2003, 2\textsuperscript{nd} Semester, LEA-pB, LEN-pB, Teaching

Students Rating: NA (practice)

1998/1999, 1\textsuperscript{st} Semester, LEEC-pB, Teaching

Students Rating: 3.09/5 (labs)

1997/1998, 1\textsuperscript{st} Semester, LEEC-pB, Teaching

Students Rating: 3.81/5 (labs)

1995/1996, 1\textsuperscript{st} Semester, LEEC-pB, Teaching

Students Rating: NA (labs)
\end{quote}

\begin{enumerate}
\def\labelenumi{\arabic{enumi}.}
\setcounter{enumi}{18}
\item
  \ul{Circuits Theory}
\end{enumerate}

2002/2003, 1\textsuperscript{st} Semestre, LEIC-pB, Teaching

Students Rating: 4.14/5 (practice)

\begin{enumerate}
\def\labelenumi{\arabic{enumi}.}
\setcounter{enumi}{19}
\item
  \ul{Signal Processing}
\end{enumerate}

2002/2003, 1\textsuperscript{st} Semester, LEA-pB, Teaching

Students Rating: 3.44/5 (practice)

\begin{enumerate}
\def\labelenumi{\arabic{enumi}.}
\setcounter{enumi}{20}
\item
  \ul{Signal Theory}
\end{enumerate}

1996/1997, 2\textsuperscript{nd} Semester, LEEC-pB, Teaching

Students Rating: 3.97/5 (practice)

\subsubsection{\texorpdfstring{Tutorial Courses (1\textsuperscript{st}
and 2\textsuperscript{nd}
Cycle)}{Tutorial Courses (1st and 2nd Cycle)}}\label{tutorial-courses-1st-and-2nd-cycle}

\begin{enumerate}
\def\labelenumi{\arabic{enumi}.}
\setcounter{enumi}{21}
\item
  \ul{2nd Cycle Integrated Project in Electrical and Computer
  Engineering}
\end{enumerate}

2025/2026, 1\textsuperscript{st} Semester, MEEC21, Teaching

2024/2025, 2\textsuperscript{nd} Semester, MEEC21, Teaching

2024/2025, 1\textsuperscript{st} Semester, MEEC21, Teaching

2023/2024, 1\textsuperscript{st} Semester, MEEC21, Teaching

2022/2023, 2\textsuperscript{nd} Semester, MEEC21, Teaching

2022/2023, 1\textsuperscript{st} Semester, MEEC21, Teaching

2021/2022, 2\textsuperscript{nd} Semester, MEEC21, Teaching

2021/2022, 1\textsuperscript{st} Semester, MEEC21, Teaching

\begin{enumerate}
\def\labelenumi{\arabic{enumi}.}
\setcounter{enumi}{22}
\item
  \ul{2nd Cycle Integrated Project in Computer Science and Engineering}
\end{enumerate}

2024/2025, 1\textsuperscript{st} Semester, MEIC-A, MEIC-T, Teaching

2023/2024, 2\textsuperscript{nd} Semester, MEIC-A, MEIC-T, Teaching

2023/2024, 1\textsuperscript{st} Semester, MEIC-A, MEIC-T, Teaching

2022/2023, 2\textsuperscript{nd} Semester, MEIC-A, MEIC-T, Teaching

\begin{quote}
2021/2022, 1\textsuperscript{st} Semester, MEIC-A, MEIC-T, Teaching

2020/2021, 2\textsuperscript{nd} Semester, MEIC-A, MEIC-T, Teaching
\end{quote}

\begin{enumerate}
\def\labelenumi{\arabic{enumi}.}
\setcounter{enumi}{23}
\item
  \ul{2nd Cycle Integrated Project in Biomedical Engineering}
\end{enumerate}

2023/2024, 1\textsuperscript{st} Semester, MEBiom21, Teaching

\begin{enumerate}
\def\labelenumi{\arabic{enumi}.}
\setcounter{enumi}{24}
\item
  \ul{Introduction to the Research in Electrical and Computer
  Engineering}
\end{enumerate}

\begin{quote}
2020/2021, 2\textsuperscript{nd} Semester, MEEC, Teaching

2020/2021, 1\textsuperscript{st} Semester, MEEC, Teaching

2019/2020, 2\textsuperscript{nd} Semester, MEEC, Teaching

2019/2020, 1\textsuperscript{st} Semester, MEEC, Teaching

2018/2019, 2\textsuperscript{nd} Semester, MEEC, Teaching

2018/2019, 1\textsuperscript{st} Semester, MEEC, Teaching

2017/2018, 2\textsuperscript{nd} Semester, MEEC, Teaching

2017/2018, 1\textsuperscript{st} Semester, MEEC, Teaching

2016/2017, 2\textsuperscript{nd} Semester, MEEC, Teaching

2016/2017, 1\textsuperscript{st} Semester, MEEC, Teaching

2015/2016, 2\textsuperscript{nd} Semester, MEEC, Teaching

2015/2016, 1\textsuperscript{st} Semester, MEEC, Teaching

2014/2015, 1\textsuperscript{st} Semester, MEEC, Teaching

2013/2014, 2\textsuperscript{nd} Semester, MEEC, Teaching

2013/2014, 1\textsuperscript{st} Semester, MEEC, Teaching

2012/2013, 1\textsuperscript{st} Semester, MEEC, Teaching
\end{quote}

\begin{enumerate}
\def\labelenumi{\arabic{enumi}.}
\setcounter{enumi}{25}
\item
  \ul{Master Project in Information and Software Engineering}
\end{enumerate}

2020/2021, 1\textsuperscript{st} Semester, MEIC-A, MEIC-T, Teaching

2018/2019, 2\textsuperscript{nd} Semester, MEIC-A, MEIC-T, Teaching

2018/2019, 1\textsuperscript{st} Semester, MEIC-A, MEIC-T, Teaching

\begin{enumerate}
\def\labelenumi{\arabic{enumi}.}
\setcounter{enumi}{26}
\item
  \ul{Project in Electronics Engineering}
\end{enumerate}

2018/2019, 1\textsuperscript{st} Semester, MEE, Teaching

\begin{enumerate}
\def\labelenumi{\arabic{enumi}.}
\setcounter{enumi}{27}
\item
  \ul{Introduction to Biomedical Engineering}
\end{enumerate}

2017/2018, 1\textsuperscript{st} Semester, MEBiom, Teaching

Students Rating: NA (project)

2016/2017, 1\textsuperscript{st} Semester, MEBiom, Teaching

Students Rating: NA (project)

2015/2016, 1\textsuperscript{st} Semester, MEBiom, Teaching

Students Rating: NA (project)

\begin{enumerate}
\def\labelenumi{\arabic{enumi}.}
\setcounter{enumi}{28}
\item
  \ul{1st Cycle Integrated Project in Biomedical Engineering}
\end{enumerate}

2023/2024, 2\textsuperscript{nd} Semester, LEBiom, Teaching

\subsection{Innovation}\label{innovation}

The following list highlights my most relevant contributions to
innovation, with emphasis in experimental and laboratory resources. Most
of the points are associated to pedagogical materials that are listed in
the next section.

\begin{enumerate}
\def\labelenumi{\arabic{enumi}.}
\item
  In 1996 I prepared, with the supervision of Prof. Luis Borges de
  Almeida a set of 10 original assignments for the laboratories of the
  Neural Networks course.
\item
  In 1998 I prepared, with the supervision of Prof. José Santos-Victor,
  a new project for the course of Computer Vision, on 3D reconstruction
  from cast shadows.
\item
  In 2002 I prepared, with the supervision of Prof. João Miranda Lemos,
  a new project for the course of Modeling, Identification and Digital
  Control, on computer-based control of a ball-in-a-beam didactic kit.
\item
  In 2003 I prepared, with the supervision of Prof. Paulo Oliveira, a
  new project for the course in Automation of Industrial Processes,
  composed of a large automation cell with 3 conveyor belts and 2 robot
  manipulators.
\item
  In 2003, I prepared, with the supervision of Prof. João Miranda Lemos,
  a Matlab/Simulink implementation of the adaptive controller MUSMAR for
  the labs of the course of Optimal and Adaptive Control.
\item
  From 2003-2007 I participated with Prof. João Miranda Lemos in several
  proposals for the improvement of the experimental teaching of several
  courses of the Scientific Area on Systems, Decision, and Control,
  through the acquisition of didactic kits (total investment exceeding
  50keur) and the preparation of laboratory projects for real-time
  control of mechanical systems. These kits have been used by several
  hundreds of students on control-based courses in Electrical and
  Computer Engineering and Aerospace Engineering since their first use
  in 2005.
\item
  In 2004, I prepared with Prof. João Miranda Lemos the first project
  based on the newly acquired didactic kits on applications to the
  control of a DC motor, for the course of Optimal and Adaptive Control.
\item
  In 2005, I adapted from C++ to Matlab/Simulink the project of the
  course Modelling, Identification and Digital Control. This adaptation
  allowed a better programming interface for the students that led them
  to progress faster in the project and a bigger diversity of
  experiments.
\item
  In 2006 I introduced new didactic kits and lab manuals for the course
  of Modelling, Identification and Digital Control, for the control of a
  flexible bar. Also, I prepared with Prof. João Miranda Lemos a new
  project for the course on State-Space Control on designing optimal
  controllers and observers for the flexible bar kit.
\item
  In 2008, with Prof. José Gaspar, I introduced, through a robotics
  middleware, new abstractions to the multitask programming component of
  the project of Distributed Real Time Control Systems.
\item
  In 2009, with Prof. José Gaspar, I introduced a more versatile
  microcontroller (Arduino) and a more accessible dynamical system (DC
  motor) in the project of Distributed Real Time Control Systems, which
  facilitated the accessibility of the students to the experimental
  equipment and the programming of the control algorithms.
\item
  In 2013 I introduced, for the first time in the Distributed Real Time
  Control Systems course decentralized control and distributed
  optimization techniques (distributed simplex), which enlarged the
  theoretical scope of the course.
\item
  In 2013 I created a new project for the course of Distributed Real
  Time Control Systems (Decentralized Control of a Distributed
  Illumination System) composed of 8 instrumented luminaires connected
  to an ethernet network, in which was possible to apply decentralized
  control algorithms for the first time.
\item
  In 2014, I created, for the first time, individual take home kits for
  the students of Distributed Real Time Control Systems. Each kit
  allowed students to develop a part of the work at home (a control node
  of the distributed control system), with more contact time with the
  experimental setup a better time management. Then, to develop the
  distributed control, students organized in groups of 2 or 3 and
  created a network using microcontroller serial communications (SPI or
  I2C) to implement the decentralized control algorithms.
\item
  In 2016, I introduced the Consensus Algorithm for Distributed
  Optimization in the course of Distributed Real Time Control Systems.
  The Consensus Algorithm with ADMM presents an alternative to the
  distributed simplex algorithm used in the previous years.
\item
  In 2019, I have introduced CAN-BUS modules to the Distributed
  Real-Time Control Systems Project. As the CAN-BUS is a control network
  widely used in industry, the introduction of these modules improves
  the practical relevance of the course project.
\item
  In 2021, I submitted a Pedagogical Innovation Project (PIP) to acquire
  microcontrollers in sufficient quantity to provide each student with a
  full kit for the Distributed Real Time Control System. Due to the
  pandemics, student had reluctance to gather physically in groups to
  develop the networked component of the project. With a full kit, each
  student had available the whole experimental setup of the project, and
  collaboration could be done remotely. In addition, were acquire more
  modern microcontrollers (Raspberry pi pico) with two cores, which made
  opportune the teaching of concurrency patterns in multi-core
  scenarios. The project was approved, with complementary funds of the
  department (DEECProj) and the scientific area of Systems, Decision and
  Control, to acquire support materials (tools, instruments) for
  students. The equipment was first used with success in the spring
  semester of 2021/2022.
\item
  In 2022, given the success of the idea of take-home labs, I submitted
  two more PIP applications for take-home kits for the projects and lab
  assignments of the courses of Control (with Prof. Paulo Oliveira and
  Prof. Rita Cunha) and Autonomous Systems (with Prof. Rodrigo Ventura
  and Prof. Alberto Vale). The projects were approved with additional
  funds from the scientific area of Systems, Decision and Control, and
  the kits were first used in the spring semester of 2022/2023.
\end{enumerate}

\subsection{Pedagogical Materials}\label{pedagogical-materials}

In this section I list the pedagogical materials of which I am the main
author and/or the main contributor. Many of these materials are
co-authored with colleagues.

\begin{enumerate}
\def\labelenumi{\arabic{enumi}.}
\item
  Distributed Real Time Control Systems

  \begin{enumerate}
  \def\labelenumii{\alph{enumii}.}
  \item
    Slides series in the following topics (evolved from a set of slides
    by Prof. Carlos Silvestre):

    \begin{enumerate}
    \def\labelenumiii{\roman{enumiii}.}
    \item
      Introduction to Distributed Real Time Control Systems
    \item
      Microcontrollers Hardware and Software
    \item
      Topics in C++
    \item
      PID Control
    \item
      Digital Control
    \item
      Real-Time Control
    \item
      Control Networks
    \item
      Concurrency
    \item
      Decentralized Control
    \item
      Distributed Algorithms
    \item
      Asynchronous Programming
    \item
      Soft skills: Reporting guidelines
    \end{enumerate}
  \item
    Project on distributed illumination control (with contributions from
    Prof. João Paulo Neto, Prof. João Pedro Gomes, Prof. José Gaspar,
    Dr. Ricardo Ribeiro) involving the following topics:

    \begin{enumerate}
    \def\labelenumiii{\roman{enumiii}.}
    \item
      Microcontroller Programming (Arduino, Raspberry Pi Pico)
    \item
      Assembly of electronic components and in small scale model of a
      distributed illumination system.
    \item
      Interface with Digital and Analog Inputs/Outputs.
    \item
      Actuation of a dimmable light source (LED) with PWM signals.
    \item
      Measuring Lux with an LDR.
    \item
      Implementing real-time control loops with timers and interrupts.
    \item
      Feedforward and feedback control with PID controllers.
    \item
      State-machines
    \item
      Multi-task programming
    \item
      Communications in buses with arbitration (CAN-BUS, I2C).
    \item
      Distributed algorithms for multi-node coordination.
    \item
      Decentralized control with competitive algorithms (Feedforward
      from accessible disturbance).
    \item
      Decentralized control with cooperative algorithms (optimization
      based -- Simplex, Consensus with ADMM).
    \item
      PC interfaces to microcontroller networks.
    \item
      Asynchronous Programming with Serial and Socket based
      communications.
    \end{enumerate}
  \item
    Project on distributed command, monitoring, and control of a DC
    motor, with a PC and an Arduino (with Prof. José Gaspar and Prof.
    Rodrigo Ventura).

    \begin{enumerate}
    \def\labelenumiii{\roman{enumiii}.}
    \item
      Physical Modelling of the DC motor plant.
    \item
      Identification of the physical parameters from step response.
    \item
      Simulation of the plant in Simulink.
    \item
      Controller design of a DC motor and implementation with an
      Arduino.
    \item
      Arduino-to-PC communications using Ethernet Shield.
    \item
      Implementation of software modules in ROS/YARP to perform
      distributed computation and communications between a Control PC
      and a monitoring PC.
    \end{enumerate}
  \end{enumerate}
\item
  Modelling, Identification, and Digital Control

  \begin{enumerate}
  \def\labelenumii{\alph{enumii}.}
  \item
    Slides series in the following topics (evolved from a set of slides
    by Prof. João Miranda Lemos):

    \begin{enumerate}
    \def\labelenumiii{\roman{enumiii}.}
    \item
      Introduction to Computer Controlled Systems
    \item
      Models of Discrete Systems
    \item
      Models of Sampled Systems
    \item
      Discretization of Continuous Controllers
    \item
      Non-parametric Identification: correlation analysis and spectral
      analysis
    \item
      Parametric Identification: AR, ARMA, ARX and ARMAX structures.
    \item
      Controller design with polynomial methods
    \item
      Linear Prediction Models
    \item
      Prediction Error Minimization and Maximum Likelihood estimation
    \item
      Minimum Variance Control
    \end{enumerate}
  \item
    Lab Assignment on Identification and Digital Control of a
    Servo-Motor with a Flexible Bar (with Prof. João Miranda Lemos and
    Prof. Margarida Silveira):

    \begin{enumerate}
    \def\labelenumiii{\roman{enumiii}.}
    \item
      Use Simulink Real-Time Workshop to interface a DC motor with
      analog position and velocity feedback.
    \item
      DC motor sensor calibration
    \item
      DC motor velocity and position Control
    \item
      Parameter Identification via Time Response
    \item
      PID control with anti-windup.
    \item
      Ziegler-Nichols controller tuning method.
    \item
      Signals for Identification: Chirp and PRBS.
    \item
      Non-parametric identification of the system: SPA and CRA
    \item
      Parametric Identification: ARX and ARMAX.
    \item
      Polynomial Control.
    \end{enumerate}
  \item
    Lab Assignment on Identification and Digital Control of the
    Ball-in-a-Beam (with Prof. João Miranda Lemos):

    \begin{enumerate}
    \def\labelenumiii{\roman{enumiii}.}
    \item
      Physical modelling and Matlab simulation.
    \item
      Parametric Identification from Step Response.
    \item
      Determination of Control
    \item
      Parametric Identification from data: ARX and ARMAX models.
    \item
      Polynomial Control with two degrees of freedom.
    \item
      Cascaded Control
    \item
      Adaptative Control
    \end{enumerate}
  \end{enumerate}
\item
  Machine Learning

  \begin{enumerate}
  \def\labelenumii{\alph{enumii}.}
  \item
    Slides series in the following topics (evolved by a set of slides
    from Prof. Luís Borges de Almeida and Prof. Fernando Mira da Silva):

    \begin{enumerate}
    \def\labelenumiii{\roman{enumiii}.}
    \item
      Introduction and Fundamentals of Machine Learning.
    \item
      Basics of Function Approximation.
    \item
      Regression vs Classification.
    \item
      Neural Networks: Adaline, MLP
    \item
      Backpropagation
    \item
      Data partitioning: Training, Validation, Testing
    \item
      Recurrent Networks
    \item
      Bayes Decision Theory
    \item
      Support Vector Machines
    \end{enumerate}
  \end{enumerate}
\item
  Computer Controlled Systems / State-Space Control / Optimal and
  Adaptive Control

  \begin{enumerate}
  \def\labelenumii{\alph{enumii}.}
  \item
    Lab Assignment on Identification and Computer Control of a Flexible
    Robot Arm Joint (with Prof. João Miranda Lemos and Prof. Margarida
    Silveira):

    \begin{enumerate}
    \def\labelenumiii{\roman{enumiii}.}
    \item
      Use Simulink Real-Time Workshop to Interface with the Dynamical
      System (send actuation commands and read sensor data).
    \item
      Calibrate system sensors.
    \item
      Methods to prepare data for identification: sampling, filtering
      and detrending.
    \item
      Non-parametric identification: Spectral Analysis from Chirp
      Signals.
    \item
      Parametric Identification: ARMAX model estimation with PRBS
      signals.
    \item
      LQ controller design.
    \item
      LQE estimator design.
    \item
      LQG controller design.
    \item
      Integral Control and Anti-Windup.
    \end{enumerate}
  \end{enumerate}
\item
  Optimal and Adaptive Control

  \begin{enumerate}
  \def\labelenumii{\alph{enumii}.}
  \item
    Lab Assignment on applications to the control of a DC motor (with
    Prof. João Miranda Lemos).

    \begin{enumerate}
    \def\labelenumiii{\roman{enumiii}.}
    \item
      Introduction to real-time computer control with Matlab/Simulink
    \item
      Proportional Control
    \item
      Sampling and Filtering
    \item
      System Identification with Step Response
    \item
      System Identification with Least Squares
    \item
      Pole Placement
    \item
      Recursive Identification
    \item
      Optimal Control and Minimum Variance Control
    \item
      Adaptive control with MUSMAR
    \end{enumerate}
  \end{enumerate}
\item
  Computer Vision

  \begin{enumerate}
  \def\labelenumii{\alph{enumii}.}
  \item
    Update of the set of slides supporting the lectures (with Prof.
    Santos-Victor).
  \item
    Lab assignment on 3D Scanner from casting shadows on objects (with
    Prof. José Santos-Victor):

    \begin{enumerate}
    \def\labelenumiii{\roman{enumiii}.}
    \item
      System calibration with checkerboard patterns through
      back-projection error minimization.
    \item
      Light source pose estimation with Least Squares
    \item
      Contour detection.
    \item
      3D reconstruction of the sequence of cast shadows.
    \end{enumerate}
  \end{enumerate}
\item
  Neural Networks

  \begin{enumerate}
  \def\labelenumii{\alph{enumii}.}
  \item
    Set of lab assignments with software ``NeuroSolutions'' (with Prof.
    Luís Borges de Almeida):

    \begin{enumerate}
    \def\labelenumiii{\roman{enumiii}.}
    \item
      Learning logic functions with ADALINE.
    \item
      Classification: Recognition of digits with MLP.
    \item
      Regression: Image Restoration with MLP.
    \item
      Recurrent Networks: Learning Linear Filters and Binary functions.
    \item
      Statistical Properties of Neural Networks: Estimation of mean,
      variance, median, percentiles and Bayesian posterior
      probabilities.
    \item
      Unsupervised and Supervised Learning: Radial Basis Functions
    \item
      Clustering and Vector Quantization
    \item
      Hopfield Networks
    \item
      Principal Component Analysis
    \end{enumerate}
  \end{enumerate}
\item
  Automation and Industrial Processes

  \begin{enumerate}
  \def\labelenumii{\alph{enumii}.}
  \item
    Lab assignment of ``Cell Automation for Quality Control'' (with
    Prof. Paulo Oliveira):

    \begin{enumerate}
    \def\labelenumiii{\roman{enumiii}.}
    \item
      Assembly of a Cell composed of Conveyor Belts (3), Robot
      Manipulators (2), proximity and defect sensors (5), push actuators
      (2).
    \item
      Modeling and Analysis of the Cell using Discrete Event Systems
    \item
      Trajectory Programming for a Robotic Manipulator
    \item
      Cell programming in GRAFCET using a PLC's
    \end{enumerate}
  \end{enumerate}
\item
  Algorithms and Data Structures

  \begin{enumerate}
  \def\labelenumii{\alph{enumii}.}
  \item
    Lab assignment on Tree data structures (with Prof. Carlos Bispo)

    \begin{enumerate}
    \def\labelenumiii{\roman{enumiii}.}
    \item
      Tree representations in file format
    \item
      In-fix and Post-Fix tree traversal
    \item
      Sweep depth and sweep breadth tree search.
    \item
      Ordering and Balancing tests.
    \end{enumerate}
  \end{enumerate}
\item
  Athens course at IST ``Introduction to Robotics'', 2007.

  \begin{enumerate}
  \def\labelenumii{\alph{enumii}.}
  \item
    Organization and production of contents
  \end{enumerate}
\item
  Best Summer Course on Intelligent Robotics Systems, module
  ``Perception'', 1996.

  \begin{enumerate}
  \def\labelenumii{\alph{enumii}.}
  \item
    Organization and production of contents (with Prof. José
    Santos-Victor)
  \end{enumerate}
\end{enumerate}
